\section*{Nomenclature}
\phantomsection
\addcontentsline{toc}{section}{Nomenclature}

% Abbreviations used in the text. Add rows as they are introduced.
\subsection*{Abbreviations}
\begin{longtable}[l]{@{}ll@{}}
AMI & Augmented multi-party interaction \\
API & Application programming interface \\
ASR  & Automatic speech recognition \\
BAK & Background intrusiveness (DNSMOS P.835 sub-score) \\
BSRNN & Band-split recurrent neural network \\
CARTSE & Li and Seki's online REAL-TSE system \cite{li2026cartse} \\
CPU & Central processing unit \\
DNSMOS & Deep noise suppression mean opinion score \\
ECAPA-TDNN & Emphasised channel attention, propagation and aggregation TDNN \\
ELBO & Evidence lower bound \\
EQ & Equalisation \\
FR & Fabrication rate \\
ICR   & Interferer content rate \\
ITU-T & ITU Telecommunication Standardization Sector \\
KL & Kullback--Leibler (divergence) \\
LAM & Large audio model \\
LCF-WER & Listener content fidelity word error rate \\
LLM & Large language model \\
LSTM  & Long short-term memory \\
LUFS & Loudness units relative to full scale \\
MLP & Multilayer perceptron \\
NLTK & Natural Language Toolkit \\
OVRL & Overall quality (DNSMOS P.835 sub-score) \\
REAL-TSE & Real-world target speaker extraction (challenge) \\
RMS & Root mean square \\
RNN & Recurrent neural network \\
RTF & Real-time factor \\
SAR & Signal-to-artefacts ratio \\
SDR & Signal-to-distortion ratio \\
SIG & Speech signal quality (DNSMOS P.835 sub-score) \\
SIR  & Signal-to-interference ratio \\
SI-SDR & Scale-invariant signal-to-distortion ratio \\
SNR & Signal-to-noise ratio \\
STFT  & Short-time Fourier transform \\
TDNN & Time-delay neural network \\
TF-Map & Time-frequency map \\
TSE   & Target speaker extraction \\
VAD & Voice activity detection \\
WER   & Word error rate \\
\end{longtable}



\subsection*{Symbols}

\subsubsection*{Signals and transcriptions}
\begin{longtable}[l]{@{}ll@{}}
$x$ & Mixture signal \\
$s$ & Target reference signal \\
$\hat{s}$ & Model output, the estimated target \\
$s_\textup{int}$ & Interfering speech signal \\
\end{longtable}

\subsubsection*{Model and objective}
\begin{longtable}[l]{@{}ll@{}}
$X$, $X_t$ & STFT of the mixture; frame $t$ of it \\
$X_e$ & STFT of the enrollment \\
$S$, $\hat{S}$ & STFT of the target; of the model output \\
$\Re(\cdot)$, $\Im(\cdot)$ & Real and imaginary parts of a complex value \\
$N$ & STFT window length \\
$\tilde{x}_t$, $\tilde{e}_j$ & Unit-length mixture frame; unit-length enrollment frame \\
$C_{t,j}$ & Cosine similarity of mixture frame $t$ and enrollment frame $j$ \\
$A_{t,j}$ & Softmax weight of enrollment frame $j$ for mixture frame $t$ \\
$\mu_t$, $u_t$ & Blended enrollment template; its unit-length direction \\
$m_t$, $\phi_t$ & Projection of $|X_t|$ onto $u_t$; the baseline cue $m_t u_t$ \\
$c_t$, $r_t$ & Similarity of $|X_t|$ to $u_t$; residual $|X_t| - m_t u_t$ \\
$M$, $R$ & Complex mask; residual spectrogram \\
$\alpha$, $s_\textup{proj}$ & Gain of the output along the target; $\alpha s$ \\
$e_\textup{int}$, $e_\textup{art}$ & Interference (other speaker and noise); artefact \\
% $d^{\perp}$, $e_d$ & $d$ with its target-like part removed; the part of $\hat{s}$ along $d^{\perp}$ \\  % REMOVED with the hail-mary loss, todo 2
$L_\textup{pres}$, $L_\textup{MR}$, $L_\textup{gain}$, $L_\textup{abs}$ & Loss terms \\
$w$, $w_\textup{m}$, $w_\textup{g}$ & Loss weights \\  % lambda removed with the hail-mary loss, todo 2
$\tau_\textup{pres}$, $\tau_\textup{abs}$ & Floors of $L_\textup{pres}$ ($10^{-3}$) and $L_\textup{abs}$ ($10^{-2}$) \\
\end{longtable}

\subsubsection*{AI usage citations}
To save space for explanations of how I used AI tools in specific parts of the project, I have made the specific declarations in the \nameref{sec:plagiarism} section which I reference using tags such as (\ref{ai:research}) and (\ref{ai:eq-gate}). 